\chapter{अग्रांत कक्षक और अभिक्रियाशीलता}\label{ch:b2:frontier-orbitals}

ताक पर छोड़ी गई साइक्लोपेंटाडाईन की बोतल धीरे-धीरे अपने द्वितय में बदल जाती है:
उसी \omterm{def:b2:frontier-orbitals:diels-alder}{डाईन} के दो अणु जुड़कर एक सेतुबद्ध द्विचक्रीय यौगिक बनाते हैं, और
संभव त्रिविम समावयवियों में से केवल एक बनता है। प्रथम वर्ष के खंड के वक्र तीरों में
कुछ भी नहीं बताता कि कौन-सा। दो कक्षक बताते हैं: एक अणु का उच्चतम भरा
कक्षक और दूसरे का निम्नतम खाली कक्षक। यह
अध्याय \cref{ch:b2:diatomic-mos} की द्वि-स्तरीय अन्योन्यक्रिया को
अभिक्रियाशीलता के साधन में बदलता है --- नाभिकरागी किस परमाणु पर आक्रमण करता है, किस ओर से, और
कोई \omterm{def:b2:frontier-orbitals:diels-alder}{डाईन} और कोई ऐल्कीन एक ही चरण में वलय में क्यों संयोजित होते हैं।

\begin{recall}
\Cref{ch:b2:diatomic-mos}: दो अन्योन्यक्रियारत कक्षक, चार-इलेक्ट्रॉन प्रतिकर्षण।
\Cref{ch:b2:huckel}: संयुग्मित निकायों के ह्युकेल स्तर और गुणांक।
प्रथम वर्ष का खंड: नाभिकरागी और इलेक्ट्रॉनरागी, वक्र तीर, गतिक
नियंत्रण, त्रिविमविशिष्ट, त्रिविमवरणात्मक और स्थानवरणात्मक अभिक्रियाएँ,
ऐल्कीनों पर इलेक्ट्रॉनरागी योग, नाभिकरागी प्रतिस्थापन।
\end{recall}

\section{दो कक्षक, फिर से}

जब दो अणु निकट आते हैं, तो उनके कक्षक अतिव्यापित होकर अन्योन्यक्रिया करते हैं। अधिकांश
अन्योन्यक्रिया दुर्बल होती है, क्योंकि कक्षक ऊर्जा में दूर होते हैं; तब द्वि-स्तरीय परिणाम का
एक विक्षोभ रूप पर्याप्त होता है।

\begin{theorem}[दो स्तरों की दुर्बल अन्योन्यक्रिया]\label{thm:b2:frontier-orbitals:perturbation}
ऊर्जाओं $\varepsilon_1 < \varepsilon_2$ वाले दो कक्षक, जो \omterm{def:b2:diatomic-mos:integrals}{अनुनाद समाकल}
$\beta$ से युग्मित हों, जहाँ $|\beta| \ll \Delta\varepsilon = \varepsilon_2 - \varepsilon_1$ (अतिव्यापन नगण्य), इन मानों पर
विस्थापित होते हैं:
\[
  \varepsilon_1 - \frac{\beta^2}{\Delta\varepsilon}, \qquad \varepsilon_2 + \frac{\beta^2}{\Delta\varepsilon}
\]
$\beta^2/\Delta\varepsilon^2$ में प्रथम कोटि तक: निचला स्तर गिरता है और ऊपरी उतना ही
उठता है, जो अंतराल जितना बड़ा उतना ही कम।
\end{theorem}

\begin{proof}
\cref{thm:b2:diatomic-mos:heteronuclear} से यथातथ स्तर $\bar\varepsilon \mp
\sqrt{\Delta\varepsilon^2/4 + \beta^2}$ हैं। $u = 4\beta^2/\Delta\varepsilon^2 \ll 1$ के साथ, $\sqrt{\Delta\varepsilon^2/4
+ \beta^2} = \frac{\Delta\varepsilon}{2}\sqrt{1 + u} \approx \frac{\Delta\varepsilon}{2}(1 + u/2) =
\frac{\Delta\varepsilon}{2} + \frac{\beta^2}{\Delta\varepsilon}$, और $\bar\varepsilon \mp \Delta\varepsilon/2 = \varepsilon_1,
\varepsilon_2$।
\end{proof}

\begin{proposition}[दो इलेक्ट्रॉन आकर्षित करते हैं, चार प्रतिकर्षित]\label{prop:b2:frontier-orbitals:four-electron}
यदि दोनों अन्योन्यक्रियारत कक्षकों में कुल दो इलेक्ट्रॉन हों, तो निकाय लगभग
$2\beta^2/\Delta\varepsilon$ से स्थायी होता है। यदि उनमें चार हों, तो निकाय अस्थायी होता है।
\end{proposition}

\begin{proof}
दो इलेक्ट्रॉन नीचे आए स्तर में जाते हैं: लाभ $2\beta^2/\Delta\varepsilon$। चार के साथ ऊपरी स्तर
भी भरता है; \cref{thm:b2:frontier-orbitals:perturbation} की कोटि तक दोनों
विस्थापन कट जाते हैं, और अतिव्यापन को रखने पर ऊपरी स्तर उतना उठता है जितना निचला
गिरता है उससे अधिक (\cref{prop:b2:diatomic-mos:asymmetry}): एक शुद्ध प्रतिकर्षण।
\end{proof}

\begin{omfigure}
\begin{tikzpicture}[font=\small]
  \begin{scope}
    \pic (a) at (0,0) {omlevel={ud}{0.7}};
    \pic (b) at (3,1.4) {omlevel={empty}{0.7}};
    \pic (lo) at (1.5,-0.5) {omlevel={ud}{0.7}};
    \pic (hi) at (1.5,1.9) {omlevel={empty}{0.7}};
    \draw[omcorr, densely dashed] (a-r) -- (lo-l) (a-r) -- (hi-l) (b-l) -- (lo-r) (b-l) -- (hi-r);
    \node[below=6pt] at (1.5,-0.6) {दो इलेक्ट्रॉन: स्थायीकृत};
    \node[left] at (a-l) {भरा};
    \node[right] at (b-r) {खाली};
  \end{scope}
  \begin{scope}[xshift=7cm]
    \pic (a) at (0,0) {omlevel={ud}{0.7}};
    \pic (b) at (3,0.6) {omlevel={ud}{0.7}};
    \pic (lo) at (1.5,-0.5) {omlevel={ud}{0.7}};
    \pic (hi) at (1.5,1.25) {omlevel={ud}{0.7}};
    \draw[omcorr, densely dashed] (a-r) -- (lo-l) (a-r) -- (hi-l) (b-l) -- (lo-r) (b-l) -- (hi-r);
    \node[below=6pt] at (1.5,-0.6) {चार इलेक्ट्रॉन: प्रतिकर्षित};
    \node[left] at (a-l) {भरा};
    \node[right] at (b-r) {भरा};
  \end{scope}
\end{tikzpicture}

{\small पड़ोसी अणुओं के दो कक्षकों की अन्योन्यक्रिया। खाली कक्षक के सामने भरा
कक्षक शुद्ध स्थायीकरण देता है; दो भरे कक्षक प्रतिकर्षित करते हैं, ऊपरी
स्तर उतना उठता है जितना निचला गिरता है उससे अधिक। यही प्रतिकर्षण अणुओं के बीच
त्रिविम बाधा का मूल है।}
\end{omfigure}

\section{अग्रांत कक्षक}

\begin{definition}[अग्रांत कक्षक]\label{def:b2:frontier-orbitals:frontier}
किसी अणु का \emph{HOMO}\index{HOMO} (उच्चतम भरा \omterm{def:b2:diatomic-mos:lcao}{आण्विक कक्षक}) और
\emph{LUMO}\index{LUMO} (निम्नतम खाली \omterm{def:b2:diatomic-mos:lcao}{आण्विक कक्षक}) उसके
\emph{अग्रांत कक्षक}\index{अग्रांत कक्षक} हैं।
\end{definition}

\begin{proposition}[प्रभावी अन्योन्यक्रिया]\label{prop:b2:frontier-orbitals:dominant}
दो संवृत-कोश अणुओं के बीच सबसे अधिक महत्त्वपूर्ण स्थायीकारी अन्योन्यक्रियाएँ
एक के \omterm{def:b2:frontier-orbitals:frontier}{HOMO} और दूसरे के \omterm{def:b2:frontier-orbitals:frontier}{LUMO} के बीच होती हैं; ऐसे दो युग्मों में से
छोटे अंतराल वाला प्रभावी होता है। नाभिकरागी अपने \omterm{def:b2:frontier-orbitals:frontier}{HOMO} से अभिक्रिया करता है,
इलेक्ट्रॉनरागी अपने \omterm{def:b2:frontier-orbitals:frontier}{LUMO} से।
\end{proposition}

\begin{proof}[तर्क]
भरा--भरा युग्म प्रतिकर्षित करते हैं (\cref{prop:b2:frontier-orbitals:four-electron}) और खाली--खाली
युग्मों में कोई इलेक्ट्रॉन नहीं होता: केवल भरा--खाली युग्म स्थायी करते हैं, हर एक $2\beta^2/\Delta\varepsilon$ से।
\omterm{def:b2:frontier-orbitals:frontier}{HOMO} उच्चतम भरा स्तर है और \omterm{def:b2:frontier-orbitals:frontier}{LUMO} निम्नतम खाली, इसलिए
\omterm{def:b2:frontier-orbitals:frontier}{HOMO} और \omterm{def:b2:frontier-orbitals:frontier}{LUMO} के युग्म का अंतराल उनमें सबसे छोटा है; प्रमेय से सबसे छोटा अंतराल
सबसे बड़ा पद देता है।
\end{proof}

\begin{definition}[आवेश नियंत्रण और कक्षक नियंत्रण]\label{def:b2:frontier-orbitals:control}
कोई अभिक्रिया \emph{आवेश नियंत्रण}\index{आवेश नियंत्रण} में है जब दोनों भागीदारों के
आवेशों के बीच आकर्षण तय करता है कि वे कहाँ अभिक्रिया करते हैं, और
\emph{कक्षक नियंत्रण}\index{कक्षक नियंत्रण} में जब \omterm{def:b2:frontier-orbitals:frontier}{अग्रांत कक्षकों} का अतिव्यापन
यह तय करता है, आक्रमण वहाँ होता है जहाँ उनके गुणांक सबसे बड़े हों।
\end{definition}

\begin{definition}[कठोर और मृदु स्पीशीज़]\label{def:b2:frontier-orbitals:hard-soft}
\emph{कठोर नाभिकरागी}\index{कठोर नाभिकरागी} या \emph{कठोर इलेक्ट्रॉनरागी}
\index{कठोर इलेक्ट्रॉनरागी} छोटा होता है, केंद्रित आवेश रखता है और उसका
\omterm{def:b2:frontier-orbitals:frontier}{HOMO} (क्रमशः \omterm{def:b2:frontier-orbitals:frontier}{LUMO}) अपने भागीदारों के कक्षक से दूर होता है: वह \omterm{def:b2:frontier-orbitals:control}{आवेश
नियंत्रण} में अभिक्रिया करता है। \emph{मृदु नाभिकरागी}\index{मृदु नाभिकरागी} या \emph{मृदु इलेक्ट्रॉनरागी}
\index{मृदु इलेक्ट्रॉनरागी} बड़ा और ध्रुवणीय होता है, ऊँचे \omterm{def:b2:frontier-orbitals:frontier}{HOMO}
(क्रमशः नीचे \omterm{def:b2:frontier-orbitals:frontier}{LUMO}) के साथ: वह \omterm{def:b2:frontier-orbitals:control}{कक्षक नियंत्रण} में अभिक्रिया करता है। कठोर वरीयता से
कठोर से अभिक्रिया करता है, मृदु मृदु से।
\end{definition}

हाइड्रॉक्साइड आयन और फ़्लुओराइड आयन \omterm{def:b2:frontier-orbitals:hard-soft}{कठोर नाभिकरागी} हैं, आयोडाइड और थायोलेट
मृदु; प्रोटॉन और कार्बधनायन \omterm{def:b2:frontier-orbitals:hard-soft}{कठोर इलेक्ट्रॉनरागी} हैं, हैलोऐल्केन का
कार्बन मृदु।

\begin{definition}[उभयदंती नाभिकरागी]\label{def:b2:frontier-orbitals:ambident}
\emph{उभयदंती नाभिकरागी}\index{उभयदंती नाभिकरागी} के दो नाभिकरागी परमाणु होते हैं
जो संयुग्मन से जुड़े हों, और जिनमें से हर एक किसी इलेक्ट्रॉनरागी पर आक्रमण कर सकता है।
\end{definition}

\omterm{def:b2:frontier-orbitals:hard-soft}{कठोर इलेक्ट्रॉनरागी} सायनाइड आयन पर नाइट्रोजन पर आक्रमण करते हैं, जहाँ ऋणात्मक आवेश अधिक है,
और \omterm{def:b2:frontier-orbitals:hard-soft}{मृदु इलेक्ट्रॉनरागी} कार्बन पर, जहाँ उसके \omterm{def:b2:frontier-orbitals:frontier}{HOMO} का गुणांक बड़ा है:
हैलोऐल्केन नाइट्राइल \ce{R-C#N} देता है। \cref{ch:b2:enolates-aldol} में मिलने वाले
ईनोलेट आयन इसी प्रकार कार्बन या ऑक्सीजन पर अभिक्रिया करते हैं।

\begin{method}[अग्रांत-कक्षक विश्लेषण]\label{met:b2:frontier-orbitals:fmo}
\begin{enumerate}
  \item नाभिकरागी (ऊँचा \omterm{def:b2:frontier-orbitals:frontier}{HOMO}) और इलेक्ट्रॉनरागी (नीचा \omterm{def:b2:frontier-orbitals:frontier}{LUMO}) पहचानिए।
  \item दोनों \omterm{def:b2:frontier-orbitals:frontier}{अग्रांत कक्षकों} को उनके गुणांकों और चिह्नों सहित खींचिए।
  \item अभिक्रिया वहाँ होती है जहाँ अतिव्यापन सबसे बड़ा हो: सबसे बड़े
        गुणांकों वाले परमाणु, एक ही चिह्न की पालियाँ आमने-सामने।
  \item निकट आने की दिशा \omterm{def:b2:frontier-orbitals:frontier}{LUMO} की आकृति का अनुसरण करती है (किसी
        \ce{C=O} तल के लंबवत, किसी \ce{C-X} आबंध के पीछे से)।
  \item \omterm{def:b2:frontier-orbitals:frontier}{HOMO} और \omterm{def:b2:frontier-orbitals:frontier}{LUMO} के बीच छोटा अंतराल तेज़ अभिक्रिया का अर्थ रखता है (अन्य बातें समान रहने पर)।
\end{enumerate}
\end{method}

आक्रमण की दिशा \omterm{def:b2:frontier-orbitals:frontier}{LUMO} से निकलती है। कार्बोनिल समूह का \omterm{def:b2:frontier-orbitals:frontier}{LUMO} उसका
$\pi^*$ कक्षक है, जो कार्बन पर बड़ा और तल के लंबवत है: नाभिकरागी
कार्बन तक तल के ऊपर या नीचे से पहुँचते हैं, \ce{C=O} अक्ष के अनुदिश नहीं।
हैलोऐल्केन का \omterm{def:b2:frontier-orbitals:frontier}{LUMO} \ce{C-X} आबंध का $\sigma^*$ कक्षक है, जिसकी कार्बन पर बड़ी पालि
X से दूर इंगित करती है: नाभिकरागी पीछे से आक्रमण करता है, और कार्बन प्रतीपित होता है
--- प्रथम वर्ष के खंड के द्वि-अणुक प्रतिस्थापन का त्रिविम-रसायन।

\section{डील्स--ऐल्डर अभिक्रिया}

\begin{definition}[समन्वित अभिक्रिया]\label{def:b2:frontier-orbitals:concerted}
\emph{समन्वित अभिक्रिया}\index{समन्वित अभिक्रिया} अपने सभी आबंध
एक ही चरण में, एक संक्रमण अवस्था से होकर, बिना किसी मध्यवर्ती के बनाती और तोड़ती है।
\end{definition}

\begin{definition}[डील्स--ऐल्डर अभिक्रिया]\label{def:b2:frontier-orbitals:diels-alder}
\emph{डील्स--ऐल्डर अभिक्रिया}\index{डील्स--ऐल्डर अभिक्रिया} अपने s-\textit{cis} संरूपण में किसी
संयुग्मित \emph{डाईन}\index{डाईन} का किसी ऐल्कीन या
ऐल्काइन, \emph{डाईनरागी}\index{डाईनरागी}, पर समन्वित योग है, जो दो
नए $\sigma$ आबंधों और एक नए $\pi$ आबंध के साथ एक षट्सदस्यी वलय बनाता है।
\end{definition}

\begin{omfigure}
\begin{tikzpicture}[font=\small, >={Stealth}, scale=1.25]
  % diene s-cis: C1 (0,1.6) C2 (0.6,2.4) C3 (1.6,2.4) C4 (2.2,1.6); dienophile below
  \draw[thick] (0,1.6) -- (0.6,2.4) -- (1.6,2.4) -- (2.2,1.6);
  \draw[thick] (0.2,1.62) -- (0.65,2.22);
  \draw[thick] (2.0,1.62) -- (1.55,2.22);
  \draw[thick] (0.2,0) -- (2.0,0); \draw[thick] (0.35,0.14) -- (1.85,0.14);
  \node[font=\scriptsize, anchor=east] at (-0.05,1.6) {1};
  \node[font=\scriptsize, anchor=south east] at (0.5,2.42) {2};
  \node[font=\scriptsize, anchor=south] at (1.6,2.42) {3};
  \node[font=\scriptsize, anchor=west] at (2.25,1.6) {4};
  % arrows: dienophile pi -> C1..C(dienophile); C1=C2 pi -> C2-C3; C3=C4 pi -> C4-C
  \draw[->, thick, omProp] (1.1,0.22) .. controls (0.9,0.9) and (0.3,0.9) .. (0.12,1.42);
  \draw[->, thick, omProp] (0.45,1.95) .. controls (0.6,2.75) and (1.0,2.8) .. (1.1,2.48);
  \draw[->, thick, omProp] (1.78,1.92) .. controls (2.5,1.6) and (2.4,0.6) .. (2.05,0.2);
  \node at (3.3,1.2) {$\longrightarrow$};
  \begin{scope}[xshift=4.2cm]
    \draw[thick] (0,1.6) -- (0.6,2.4) -- (1.6,2.4) -- (2.2,1.6) -- (2.0,0) -- (0.2,0) -- cycle;
    \draw[thick] (0.72,2.26) -- (1.48,2.26);
  \end{scope}
\end{tikzpicture}

{\small ब्यूटाडाईन और एथीन की \omterm{def:b2:frontier-orbitals:diels-alder}{डील्स--ऐल्डर अभिक्रिया}, वक्र तीरों से लिखी गई:
इलेक्ट्रॉनों के तीन युग्म एक साथ, एक चक्रीय, समन्वित चरण में चलते हैं; \omterm{def:b2:frontier-orbitals:diels-alder}{डाईन} के सिरों पर दो $\sigma$ आबंध
बनते हैं, और $\pi$ आबंध उसके केंद्र पर खिसक जाता है।}
\end{omfigure}

प्रभावी अन्योन्यक्रिया \omterm{def:b2:frontier-orbitals:diels-alder}{डाईन} के \omterm{def:b2:frontier-orbitals:frontier}{HOMO}, $\psi_2$, और
\omterm{def:b2:frontier-orbitals:diels-alder}{डाईनरागी} के \omterm{def:b2:frontier-orbitals:frontier}{LUMO}, $\pi^*$, के बीच है। \omterm{def:b2:frontier-orbitals:diels-alder}{डाईन} के सिरों पर और \omterm{def:b2:frontier-orbitals:diels-alder}{डाईनरागी} के दोनों कार्बनों पर
उनकी पालियों के चिह्न मेल खाते हैं जब \omterm{def:b2:frontier-orbitals:diels-alder}{डाईनरागी} \omterm{def:b2:frontier-orbitals:diels-alder}{डाईन} के फलक की ओर आता है:
दोनों नए आबंध एक साथ, हर भागीदार के एक ही फलक पर, बन सकते हैं।

\begin{omfigure}
\begin{tikzpicture}[font=\small]
  \foreach \x/\c in {0/0.602, 1/0.372, 2/-0.372, 3/-0.602} {
    \pic at (\x,2) {omp={1.5*\c}{90}};
  }
  \draw[black!45] (-0.2,2) -- (3.2,2);
  \node[anchor=west] at (3.6,2) {डाईन का HOMO ($\psi_2$)};
  \pic at (0,0) {omp={-0.85}{90}}; \pic at (3,0) {omp={0.85}{90}};
  \draw[black!45] (-0.2,0) -- (3.2,0);
  \node[anchor=west] at (3.6,0) {डाईनरागी का LUMO ($\pi^*$)};
  \draw[omProp, thick, densely dashed] (0,1.45) -- (0,0.6);
  \draw[omProp, thick, densely dashed] (3,1.45) -- (3,0.6);
  \node[font=\scriptsize, omProp, anchor=east] at (-0.15,1.0) {समकला में};
  \node[font=\scriptsize, omProp, anchor=west] at (3.15,1.0) {समकला में};
\end{tikzpicture}

{\small ब्यूटाडाईन और एथीन के \omterm{def:b2:frontier-orbitals:frontier}{अग्रांत कक्षक}, पालियाँ उनके
ह्युकेल गुणांकों के अनुपात में। आमने-सामने होने पर \omterm{def:b2:frontier-orbitals:diels-alder}{डाईन} के \omterm{def:b2:frontier-orbitals:frontier}{HOMO} की सिरे वाली पालियों और
\omterm{def:b2:frontier-orbitals:diels-alder}{डाईनरागी} के \omterm{def:b2:frontier-orbitals:frontier}{LUMO} की पालियों के चिह्न दोनों सिरों पर समान हैं: दोनों
$\sigma$ आबंध साथ-साथ, हर भागीदार के एक ही फलक पर, बन सकते हैं।}
\end{omfigure}

\begin{proposition}[त्रिविमविशिष्टता]\label{prop:b2:frontier-orbitals:da-stereospecific}
\omterm{def:b2:frontier-orbitals:diels-alder}{डील्स--ऐल्डर अभिक्रिया} त्रिविमविशिष्ट है: \omterm{def:b2:frontier-orbitals:diels-alder}{डाईनरागी} पर cis प्रतिस्थापी
उत्पाद में cis रहते हैं, trans प्रतिस्थापी trans ही रहते हैं।
\end{proposition}

\begin{proof}
अभिक्रिया समन्वित है और दोनों नए आबंध \omterm{def:b2:frontier-orbitals:diels-alder}{डाईनरागी} के एक ही फलक पर बनते
हैं: उसके दोनों कार्बन कभी एक-दूसरे के सापेक्ष नहीं घूमते, और उनके प्रतिस्थापियों की
आपेक्षिक स्थितियाँ वलय में चली आती हैं।
\end{proof}

\begin{proposition}[स्थानवरणात्मकता]\label{prop:b2:frontier-orbitals:da-regio}
असममित भागीदारों के साथ मुख्य उत्पाद \omterm{def:b2:frontier-orbitals:diels-alder}{डाईन} के उस परमाणु को, जिसका \omterm{def:b2:frontier-orbitals:frontier}{HOMO} में
गुणांक सबसे बड़ा है, \omterm{def:b2:frontier-orbitals:diels-alder}{डाईनरागी} के उस परमाणु से जोड़ता है जिसका \omterm{def:b2:frontier-orbitals:frontier}{LUMO} में
गुणांक सबसे बड़ा है।
\end{proposition}

\begin{proof}[तर्क]
संक्रमण अवस्था में दोनों नए आबंध अभी समान रूप से नहीं बने होते;
स्थायीकरण $2\beta^2/\Delta\varepsilon$ तब बड़ा होता है जब बड़े गुणांक अतिव्यापित हों ($\beta$
अतिव्यापित परमाणुओं के गुणांकों के गुणनफल के साथ बढ़ता है), इसलिए वह
अभिविन्यास कम ऊर्जा पर प्राप्त होता है।
\end{proof}

\begin{omfigure}
\begin{tikzpicture}
\begin{axis}[ybar, width=0.47\linewidth, height=4.6cm, ymin=-0.8, ymax=0.8,
  xtick={0,1,2,3,4}, xticklabels={O,C1,C2,C3,C4}, axis x line=bottom, axis y line=left,
  extra y ticks={0}, extra y tick style={grid=major, grid style={black!50}}, enlarge x limits=0.15,
  tick label style={font=\scriptsize}, ylabel={गुणांक}, label style={font=\small},
  title={\small 1-मेथॉक्सीब्यूटाडाईन का HOMO}, bar width=8pt]
\addplot[fill=omDef!50, draw=omDef] table[x=atom, y=c] {figdata/out/bachelor-2-huckel-c-diene.dat};
\end{axis}
\end{tikzpicture}\hfill
\begin{tikzpicture}
\begin{axis}[ybar, width=0.47\linewidth, height=4.6cm, ymin=-0.8, ymax=0.8,
  xtick={1,2,3,4}, xticklabels={C$\beta$,C$\alpha$,C(O),O}, axis x line=bottom, axis y line=left,
  extra y ticks={0}, extra y tick style={grid=major, grid style={black!50}}, enlarge x limits=0.15,
  tick label style={font=\scriptsize}, ylabel={गुणांक}, label style={font=\small},
  title={\small प्रोपीनल का LUMO}, bar width=8pt]
\addplot[fill=omThm!45, draw=omThm] table[x=atom, y=c] {figdata/out/bachelor-2-huckel-c-dienophile.dat};
\end{axis}
\end{tikzpicture}

{\small विषम-परमाणु प्राचलों वाला ह्युकेल मॉडल (ईथर ऑक्सीजन के लिए $\alpha_{\ce{O}} = \alpha + 2\beta$,
$\beta_{\ce{CO}} = 0.8\beta$; कार्बोनिल ऑक्सीजन के लिए $\alpha_{\ce{O}} = \alpha + \beta$, $\beta_{\ce{CO}} =
\beta$; एक मॉडल)। \omterm{def:b2:frontier-orbitals:diels-alder}{डाईन} का \omterm{def:b2:frontier-orbitals:frontier}{HOMO} C4 पर, मेथॉक्सी समूह से दूर वाले सिरे पर, सबसे बड़ा है
(0.60, जबकि 0.51); \omterm{def:b2:frontier-orbitals:diels-alder}{डाईनरागी} का \omterm{def:b2:frontier-orbitals:frontier}{LUMO}
$\beta$ कार्बन पर सबसे बड़ा है (0.66)। वे आपस में आबंधित होते हैं: मेथॉक्सी और ऐल्डिहाइड समूह
वलय के पड़ोसी कार्बनों पर पहुँचते हैं।}
\end{omfigure}

दाता समूह \omterm{def:b2:frontier-orbitals:diels-alder}{डाईन} के \omterm{def:b2:frontier-orbitals:frontier}{HOMO} को ऊपर उठाता है ($m = 0.48$, न कि 0.62) और ग्राही समूह
\omterm{def:b2:frontier-orbitals:diels-alder}{डाईनरागी} के \omterm{def:b2:frontier-orbitals:frontier}{LUMO} को नीचे लाता है ($m = -0.35$, न कि $-1$): अंतराल संकरा होता है, और ऐसे
युग्म एथीन के साथ ब्यूटाडाईन की तुलना में कहीं तेज़ी से अभिक्रिया करते हैं। इलेक्ट्रॉन-समृद्ध \omterm{def:b2:frontier-orbitals:diels-alder}{डाईन} और
इलेक्ट्रॉन-निर्धन \omterm{def:b2:frontier-orbitals:diels-alder}{डाईनरागी} विशिष्ट संयोजन है।

\begin{definition}[endo और exo योगज]\label{def:b2:frontier-orbitals:endo}
जब कोई चक्रीय \omterm{def:b2:frontier-orbitals:diels-alder}{डाईन} असंतृप्त प्रतिस्थापी वाले \omterm{def:b2:frontier-orbitals:diels-alder}{डाईनरागी} से योग करता है, तो
\emph{endo योगज}\index{endo योगज} में वह प्रतिस्थापी नए
बने द्वि-आबंध की ओर इंगित करता है (संक्रमण अवस्था में \omterm{def:b2:frontier-orbitals:diels-alder}{डाईन} के नीचे), और \emph{exo
योगज}\index{exo योगज} में उससे दूर। \emph{endo नियम}\index{endo नियम} कहता है
कि endo योगज तेज़ी से बनता है।
\end{definition}

\begin{proposition}[endo नियम]\label{prop:b2:frontier-orbitals:endo}
गतिक नियंत्रण में \omterm{def:b2:frontier-orbitals:endo}{endo योगज} मुख्य उत्पाद है, यद्यपि \omterm{def:b2:frontier-orbitals:endo}{exo योगज}
प्रायः अधिक स्थायी होता है।
\end{proposition}

\admitted

सामान्य व्याख्या endo उपागम में प्रतिस्थापी के $\pi$ निकाय की पालियों और \omterm{def:b2:frontier-orbitals:diels-alder}{डाईन} के
\omterm{def:b2:frontier-orbitals:frontier}{HOMO} के केंद्रीय परमाणुओं की पालियों के बीच एक द्वितीयक अतिव्यापन है; यह
कोई आबंध बनाए बिना endo संक्रमण अवस्था को नीचे लाता है। नियम
अनुभवजन्य है और इसके अपवाद हैं; तृतीय वर्ष का खंड इन अभिक्रियाओं का उनके कक्षकों की
सममिति के साथ विवेचन करता है।

\begin{omfigure}
\begin{tikzpicture}[font=\small]
  \foreach \xs/\mode in {0/endo, 5/exo} {
    \begin{scope}[xshift=\xs cm]
      % cyclopentadiene seen from the side: C1, C2=C3, C4, CH2 bridge on top
      \draw[thick] (0,2) -- (0.6,2.5) -- (1.8,2.5) -- (2.4,2);
      \draw[thick] (0.75,2.38) -- (1.65,2.38);
      \draw[thick] (0,2) -- (1.2,3.2) -- (2.4,2);
      \node[font=\scriptsize, anchor=south] at (1.2,3.2) {\ce{CH2}};
      % dienophile
      \draw[thick] (0.2,0) -- (2.2,0); \draw[thick] (0.35,0.1) -- (2.05,0.1);
      \draw[omProp, densely dashed, thick] (0,1.9) -- (0.2,0.15);
      \draw[omProp, densely dashed, thick] (2.4,1.9) -- (2.2,0.15);
      \node[anchor=north] at (1.2,-1.3) {\mode};
    \end{scope}
  }
  % endo: carbonyls tucked under the diene
  \draw[thick] (0.2,0) -- (0.75,0.95) (2.2,0) -- (1.65,0.95);
  \node[font=\scriptsize, fill=white, inner sep=1pt] at (0.75,1.1) {C=O};
  \node[font=\scriptsize, fill=white, inner sep=1pt] at (1.65,1.1) {C=O};
  \draw[omThm, densely dotted, thick] (0.75,1.25) -- (0.65,2.35);
  \draw[omThm, densely dotted, thick] (1.65,1.25) -- (1.75,2.35);
  % exo: carbonyls pointing away
  \draw[thick] (5.2,0) -- (4.95,-0.85) (7.2,0) -- (7.45,-0.85);
  \node[font=\scriptsize] at (4.95,-1.0) {C=O};
  \node[font=\scriptsize] at (7.45,-1.0) {C=O};
\end{tikzpicture}

{\small साइक्लोपेंटाडाईन पर मैलेइक ऐनहाइड्राइड के endo और exo उपागम, व्यवस्था-रूप।
बनते आबंध खंडित हैं। endo उपागम में कार्बोनिल समूह \omterm{def:b2:frontier-orbitals:diels-alder}{डाईन} के
नीचे होते हैं, और उनके $\pi$ कक्षक उसके केंद्रीय कार्बनों से अतिव्यापित होते हैं (बिंदुकित, द्वितीयक
अतिव्यापन); exo उपागम में वे दूर इंगित करते हैं।}
\end{omfigure}

\begin{method}[डील्स--ऐल्डर उत्पाद खींचना]\label{met:b2:frontier-orbitals:diels-alder}
\begin{enumerate}
  \item \omterm{def:b2:frontier-orbitals:diels-alder}{डाईन} को उसके s-\textit{cis} संरूपण में रखिए, \omterm{def:b2:frontier-orbitals:diels-alder}{डाईनरागी} उसके सिरों के सामने।
  \item \omterm{def:b2:frontier-orbitals:diels-alder}{डाईन} के सिरों से \omterm{def:b2:frontier-orbitals:diels-alder}{डाईनरागी} के कार्बनों तक दो नए $\sigma$ आबंध खींचिए,
        और \omterm{def:b2:frontier-orbitals:diels-alder}{डाईन} के C2 और C3 के बीच नया $\pi$ आबंध।
  \item \omterm{def:b2:frontier-orbitals:diels-alder}{डाईनरागी} के प्रतिस्थापियों को उसी ओर रखिए (cis, cis ही रहता है)।
  \item असममित भागीदारों के लिए सबसे बड़े गुणांकों को जोड़िए (1- या 2-प्रतिस्थापित \omterm{def:b2:frontier-orbitals:diels-alder}{डाईनों}
        के लिए ``ऑर्थो'' या ``पैरा'' उत्पाद)।
  \item चक्रीय \omterm{def:b2:frontier-orbitals:diels-alder}{डाईन} के साथ असंतृप्त प्रतिस्थापी को endo रखिए।
\end{enumerate}
\end{method}

\begin{inthelab}[डाइसाइक्लोपेंटाडाईन का भंजन]
साइक्लोपेंटाडाईन उसके द्वितय के रूप में बेची जाती है और उपयोग से ठीक पहले पुनर्जनित की जाती है: द्वितय को
प्रभाजी स्तंभ लगे फ़्लास्क में उसके क्वथनांक (लगभग \qty{170}{\degreeCelsius}) तक
गर्म किया जाता है, जहाँ वह वापस एकलक में टूट जाता है (एक प्रतीप \omterm{def:b2:frontier-orbitals:diels-alder}{डील्स--ऐल्डर
अभिक्रिया}), जो \qty{41}{\degreeCelsius} पर आसवित होकर बर्फ़ में ठंडे किए गए ग्राही में
संग्रहित होता है। कमरे के ताप पर एकलक कुछ ही घंटों में फिर द्वितय बन जाता है और उसे
तुरंत प्रयुक्त किया जाता है। दोनों यौगिक ज्वलनशील और हानिकारक हैं; काम धूम-हुड में किया जाता है।
\end{inthelab}

\begin{safety}{\ghs[5mm]{GHS02}\,\ghs[5mm]{GHS06}\,\ghs[5mm]{GHS08}}
साइक्लोपेंटाडाईन और उसका द्वितय अत्यंत ज्वलनशील और साँस लेने पर विषैले हैं; मैलेइक
ऐनहाइड्राइड संक्षारक और श्वसन संवेदीकारक है। इन्हें धूम-हुड में,
दस्तानों और नेत्र-सुरक्षा के साथ, ज्वालाओं से दूर, सँभाला जाता है।
% ledger: ghs:cyclopentadiene, ghs:dicyclopentadiene, ghs:maleic-anhydride, bp:cyclopentadiene, bp:dicyclopentadiene
\end{safety}

\section{अभ्यास}

\begin{exercise}[$\star$]\label{exo:b2:frontier-orbitals:1}
एथीन, ऐलिल ऋणायन
और ब्यूटाडाईन का \omterm{def:b2:frontier-orbitals:frontier}{HOMO} और \omterm{def:b2:frontier-orbitals:frontier}{LUMO} (ह्युकेल इकाइयों में ऊर्जाएँ) दीजिए।
\end{exercise}

\begin{exercise}[$\star$]\label{exo:b2:frontier-orbitals:2}
कठोर या मृदु के रूप में वर्गीकृत कीजिए: \ce{OH-}, \ce{I-}, \ce{H+}, एक थायोलेट \ce{RS-}, आयडोमेथेन
का कार्बन, \ce{F-}।
\end{exercise}

\begin{exercise}[$\star$]\label{exo:b2:frontier-orbitals:3}
इनमें से कौन-से \omterm{def:b2:frontier-orbitals:diels-alder}{डाईन} \omterm{def:b2:frontier-orbitals:diels-alder}{डील्स--ऐल्डर अभिक्रिया} में भाग ले सकते हैं: ब्यूटा-1,3-डाईन,
साइक्लोपेंटाडाईन, पेंटा-1,4-डाईन, अपने s-\textit{cis} रूप में (2\textit{Z},4\textit{Z})-हेक्सा-2,4-डाईन?
\end{exercise}

\begin{exercise}[$\star$]\label{exo:b2:frontier-orbitals:4}
प्रोपीननाइट्राइल \ce{CH2=CH-C#N} के साथ ब्यूटाडाईन का उत्पाद खींचिए।
\end{exercise}

\begin{exercise}[$\star\star$]\label{exo:b2:frontier-orbitals:5}
$-10$ और \qty{-2}{eV} पर दो कक्षक $\beta = \qty{-1.0}{eV}$ के साथ अन्योन्यक्रिया करते हैं (अभ्यास के आँकड़े)।
दो इलेक्ट्रॉनों का विक्षोभ स्थायीकरण परिकलित कीजिए, फिर यथातथ स्थायीकरण।
\end{exercise}

\begin{exercise}[$\star\star$]\label{exo:b2:frontier-orbitals:6}
सायनाइड आयन आयडोमेथेन से अभिक्रिया करके एथेननाइट्राइल \ce{CH3-C#N} देता है, परंतु
प्रबल आयनकारी विलायक में किसी कार्बधनायन के साथ वह कुछ आइसोनाइट्राइल \ce{R-N#C} देता है। समझाइए।
\end{exercise}

\begin{exercise}[$\star\star$]\label{exo:b2:frontier-orbitals:7}
अध्याय के गुणांकों से प्रोपीनल के साथ 1-मेथॉक्सीब्यूटाडाईन के मुख्य उत्पाद का
पूर्वानुमान कीजिए।
\end{exercise}

\begin{exercise}[$\star\star$]\label{exo:b2:frontier-orbitals:8}
मेथिल प्रोपीनोएट के साथ साइक्लोपेंटाडाईन का \omterm{def:b2:frontier-orbitals:endo}{endo योगज} खींचिए।
\end{exercise}

\begin{exercise}[$\star\star$]\label{exo:b2:frontier-orbitals:9}
मैलेइक ऐनहाइड्राइड कमरे के ताप पर साइक्लोपेंटाडाईन से अभिक्रिया क्यों करता है, जबकि एथीन
को दाब के अंतर्गत गर्म करना पड़ता है?
\end{exercise}

\begin{exercise}[$\star\star\star$]\label{exo:b2:frontier-orbitals:10}
साइक्लोपेंटाडाईन डाइमेथिल मैलिएट (cis डाइएस्टर) और डाइमेथिल फ़्यूमेरेट
(trans डाइएस्टर) से अभिक्रिया करती है। दोनों उत्पाद खींचिए और बताइए कि हर एक कितने त्रिविम समावयवी देता है।
\end{exercise}

\begin{exercise}[$\star\star\star$]\label{exo:b2:frontier-orbitals:11}
\omterm{def:b2:frontier-orbitals:diels-alder}{डील्स--ऐल्डर अभिक्रिया} के लिए $\Delta_r H^\circ < 0$ और $\Delta_r S^\circ < 0$ हैं। चिह्नों को
समझाइए, और यह भी कि उच्च ताप पर प्रतीप अभिक्रिया अनुकूल क्यों हो जाती है।
\end{exercise}

\begin{exercise}[$\star\star\star$]\label{exo:b2:frontier-orbitals:12}
हाइड्रैज़ीन \ce{H2N-NH2} में हर नाइट्रोजन एक एकाकी युग्म रखता है। चार-इलेक्ट्रॉन
अन्योन्यक्रिया से समझाइए कि अणु उस संरूपण से क्यों बचता है जिसमें दोनों एकाकी
युग्म समांतर हों।
\end{exercise}

\section{समस्या: साइक्लोपेंटाडाईन का भंजन}

\begin{problem}[{सप्ताहांत समस्या --- डील्स--ऐल्डर अभिक्रिया के रूप में
साइक्लोपेंटाडाईन का द्वितयन, भंजन की ऊष्मागतिकी,
साइक्लोपेंटाडाईन और मैलेइक ऐनहाइड्राइड के अग्रांत कक्षक, और उनके योगज का त्रिविम-रसायन}]
\label{pb:b2:frontier-orbitals:1}
आँकड़े: क्वथनांक, साइक्लोपेंटाडाईन \qty{41}{\degreeCelsius}, डाइसाइक्लोपेंटाडाईन लगभग
\qty{170}{\degreeCelsius}। ह्युकेल मॉडल ऊर्जाएँ (एक मॉडल): साइक्लोपेंटाडाईन, एक
ब्यूटाडाईन इकाई के रूप में, \omterm{def:b2:frontier-orbitals:frontier}{HOMO} $\alpha + 0.62\beta$ और \omterm{def:b2:frontier-orbitals:frontier}{LUMO} $\alpha - 0.62\beta$; मैलेइक ऐनहाइड्राइड, \omterm{def:b2:frontier-orbitals:frontier}{LUMO}
लगभग $\alpha - 0.35\beta$।
% ledger: bp:cyclopentadiene, bp:dicyclopentadiene

\medskip
\noindent\textbf{भाग I --- द्वितय।}
\begin{enumerate}
  \item साइक्लोपेंटाडाईन खींचिए और दिखाइए कि उसकी \omterm{def:b2:frontier-orbitals:diels-alder}{डाईन} इकाई s-\textit{cis} में बँधी है।
  \item द्वितयन में एक अणु \omterm{def:b2:frontier-orbitals:diels-alder}{डाईन} है। दूसरा क्या है?
  \item द्वितय खींचिए और नए आबंधों को क्रमांकित कीजिए।
  \item कमरे के ताप पर कौन-सा योगज, endo या exo, बनता है, और क्यों?
  \item क्या द्वितयन त्रिविमविशिष्ट है? क्रियाविधि से समझाइए।
  \item द्वितयन को उत्प्रेरक की आवश्यकता क्यों नहीं होती?
\end{enumerate}

\medskip
\noindent\textbf{भाग II --- भंजन की ऊष्मागतिकी।}
\begin{enumerate}[resume]
  \item बने और टूटे आबंधों से द्वितयन के $\Delta_r H^\circ$ का चिह्न
        दीजिए।
  \item $\Delta_r S^\circ$ का चिह्न दीजिए।
  \item $\Delta_r G^\circ(T)$ लिखिए और दिखाइए कि यह किसी ताप $T_i$ पर चिह्न बदलता है।
  \item कमरे के ताप पर द्वितय स्थायी क्यों है, और गर्म होने पर एकलक अनुकूल क्यों?
  \item भंजन व्यवस्था में एकलक को बनते ही आसवित करना अभिक्रिया को आगे क्यों
        चलाता है?
  \item एकलक को ठंडा क्यों रखना चाहिए और जल्दी क्यों प्रयुक्त करना चाहिए?
\end{enumerate}

\medskip
\noindent\textbf{भाग III --- \omterm{def:b2:frontier-orbitals:frontier}{अग्रांत कक्षक}।}
\begin{enumerate}[resume]
  \item स्वयं से अभिक्रिया करती साइक्लोपेंटाडाईन के लिए \omterm{def:b2:frontier-orbitals:frontier}{HOMO} और \omterm{def:b2:frontier-orbitals:frontier}{LUMO} के बीच अंतराल
        $|\beta|$ की इकाइयों में परिकलित कीजिए।
  \item मैलेइक ऐनहाइड्राइड के साथ साइक्लोपेंटाडाईन के लिए \omterm{def:b2:frontier-orbitals:diels-alder}{डाईन} के
        \omterm{def:b2:frontier-orbitals:frontier}{HOMO} और ऐनहाइड्राइड के \omterm{def:b2:frontier-orbitals:frontier}{LUMO} के बीच अंतराल परिकलित कीजिए।
  \item कौन-सा युग्म तेज़ी से अभिक्रिया करता है? समझाइए।
  \item मैलेइक ऐनहाइड्राइड के दोनों कार्बोनिल समूह उसके \omterm{def:b2:frontier-orbitals:frontier}{LUMO} को नीचे क्यों लाते हैं?
  \item मैलेइक ऐनहाइड्राइड का कौन-सा कक्षक endo उपागम का द्वितीयक अतिव्यापन
        देता है?
\end{enumerate}

\medskip
\noindent\textbf{भाग IV --- मैलेइक ऐनहाइड्राइड के साथ योगज।}
\begin{enumerate}[resume]
  \item \omterm{def:b2:frontier-orbitals:endo}{endo योगज} खींचिए।
  \item उसके त्रिविमजनक केंद्र चिह्नित कीजिए।
  \item क्या योगज किरैल है? (दर्पण तल खोजिए।)
  \item भूतपूर्व \omterm{def:b2:frontier-orbitals:diels-alder}{डाईनरागी} के दोनों हाइड्रोजन वलय की एक ही ओर क्यों
        पहुँचते हैं?
  \item \omterm{def:b2:frontier-orbitals:endo}{exo योगज} कैसा दिखेगा?
  \item क्या endo और \omterm{def:b2:frontier-orbitals:endo}{exo योगज} कमरे के ताप पर एक-दूसरे में बदल सकते हैं? क्यों?
  \item ऐनहाइड्राइड जल से कैसे अभिक्रिया करेगा, और कौन-सा उत्पाद बनेगा?
  \item साइक्लोपेंटाडाईन और मैलेइक ऐनहाइड्राइड के \omterm{def:b2:frontier-orbitals:endo}{endo योगज} के त्रिविमजनक केंद्रों की
        संख्या बताइए।
\end{enumerate}
\end{problem}
